\chapter{LSTM-Autoencoder Model Architecture}
\label{ch:model_architecture}

\section{Recurrent Neural Networks and Their Limitations}
\label{sec:rnn}

Recurrent Neural Networks (RNNs) are a foundational class of neural architectures
designed to process sequential data \cite{rumelhart1986learning, goodfellow2016deep}.
Unlike feedforward networks, RNNs introduce recurrent connections that allow
information to persist across time steps. At each step $t$, the network combines the
current input $x_t$ with the previous hidden state $h_{t-1}$ to compute:

\begin{equation}
    h_t = \sigma(W_{xh} x_t + W_{hh} h_{t-1} + b_h)
    \label{eq:rnn_hidden}
\end{equation}

\noindent where $W_{xh}$ and $W_{hh}$ are learnable weight matrices, $b_h$ is a bias
vector, and $\sigma$ is a nonlinear activation (typically $\tanh$).

Despite their expressiveness, vanilla RNNs suffer from well-known training instabilities
when learning long-range dependencies \cite{bengio1994learning, pascanu2013difficulty}.
During Backpropagation Through Time (BPTT), repeated multiplication of weight matrices
across many steps causes gradients to either vanish exponentially or grow without bound.
The vanishing gradient problem is particularly harmful, as the error signal decays
before it can influence earlier states. Gradient clipping \cite{pascanu2013difficulty}
mitigates exploding gradients, but the vanishing gradient problem fundamentally
motivates the use of gated architectures such as Long Short-Term Memory networks.

\section{Long Short-Term Memory (LSTM)}
\label{sec:lstm}

\subsection{LSTM Architecture}
\label{subsec:lstm_architecture}

The LSTM network, introduced by Hochreiter and Schmidhuber
\cite{hochreiter1997long, graves2012supervised}, addresses the vanishing gradient
problem through a gating mechanism that regulates information flow within the recurrent
cell. An LSTM cell maintains a cell state $C_t$ (long-term memory) and a hidden state
$h_t$ (working memory), governed by three multiplicative gates:

\begin{itemize}
    \item \textbf{Forget Gate} ($f_t$): Determines what to discard from the previous
    cell state:
    \begin{equation}
        f_t = \sigma(W_f \cdot [h_{t-1}, x_t] + b_f)
        \label{eq:forget_gate}
    \end{equation}

    \item \textbf{Input Gate} ($i_t$): Controls what new information is written, using
    a candidate cell state $\tilde{C}_t$:
    \begin{equation}
        i_t = \sigma(W_i \cdot [h_{t-1}, x_t] + b_i)
        \label{eq:input_gate}
    \end{equation}
    \begin{equation}
        \tilde{C}_t = \tanh(W_C \cdot [h_{t-1}, x_t] + b_C)
        \label{eq:candidate_cell}
    \end{equation}

    \item \textbf{Output Gate} ($o_t$): Regulates which parts of the cell state are
    exposed as the hidden state:
    \begin{equation}
        o_t = \sigma(W_o \cdot [h_{t-1}, x_t] + b_o)
        \label{eq:output_gate}
    \end{equation}
\end{itemize}

\noindent The cell and hidden states are then updated as:

\begin{equation}
    C_t = f_t \odot C_{t-1} + i_t \odot \tilde{C}_t
    \label{eq:cell_state}
\end{equation}

\begin{equation}
    h_t = o_t \odot \tanh(C_t)
    \label{eq:hidden_state}
\end{equation}

\noindent where $\sigma$ is the sigmoid function and $\odot$ denotes element-wise
multiplication. The additive cell-state update provides a gradient highway that
preserves information across long time horizons \cite{gers2000learning}, making LSTMs
significantly more stable than vanilla RNNs for learning long-range dependencies.

\subsection{LSTM for Anomaly Detection}
\label{subsec:lstm_anomaly}

In time-series anomaly detection, LSTM-based models are trained on data representative
of normal system behavior \cite{malhotra2015lstm, hundman2018detecting,
ergen2017unsupervised}. The model learns temporal regularities and multivariate
relationships present in nominal operational patterns. During inference, sequences that
deviate from the learned representation of normality typically produce higher
reconstruction or prediction errors, enabling unsupervised fault detection without
requiring labeled anomaly examples.

\section{Autoencoders}
\label{sec:autoencoders}

\subsection{Autoencoder Architecture}
\label{subsec:ae_architecture}

An autoencoder is an unsupervised neural network that learns to encode its input into a
lower-dimensional latent representation and subsequently reconstruct the original input
\cite{hinton2006reducing, goodfellow2016deep}. The architecture consists of two
complementary sub-networks:

\begin{itemize}
    \item \textbf{Encoder} $f_\theta$: Maps the input $x \in \mathbb{R}^d$ to a
    latent vector $z \in \mathbb{R}^k$, where $k < d$.

    \item \textbf{Decoder} $g_\phi$: Reconstructs the input from the latent
    representation, producing $\hat{x} = g_\phi(z) \approx x$.
\end{itemize}

The model is trained end-to-end to minimize reconstruction error, typically using Mean
Squared Error (MSE):

\begin{equation}
    \mathcal{L}_{\text{MSE}} = \frac{1}{N} \sum_{i=1}^{N} \|x_i - \hat{x}_i\|^2
    \label{eq:mse_loss}
\end{equation}

In this expression, $N$ is the number of training examples in the batch,
$x_i$ is the input vector or flattened window, and $\hat{x}_i$ is the
corresponding reconstruction. For a telemetry window
$X \in \mathbb{R}^{T \times d}$, the norm expands over all time steps and
features, so reconstruction quality depends on both temporal shape and
cross-feature relationships. The bottleneck imposed by the latent dimension
$k$ forces the model to learn a compressed representation of the input
distribution rather than merely replicating the raw input.

\subsection{Autoencoders for Anomaly Detection}
\label{subsec:ae_why}

Autoencoders are widely used for anomaly detection because they learn
representations of normal data without explicit anomaly labels
\cite{chalapathy2019deep, pang2021deep}: (1)~unsupervised training on normal
behavior, where anomalous examples are rare; (2)~multivariate dependency
modeling through compressed latent representations; and (3)~reconstruction
error as a natural anomaly score, since out-of-distribution inputs produce
elevated errors. Compared to PCA---a linear analog---neural autoencoders
capture nonlinear feature interactions \cite{jolliffe2016pca} and extend
naturally to sequential variants such as the LSTM-AE.

\section{LSTM-Autoencoder Architecture}
\label{sec:lstm_ae}

\subsection{Combined Architecture}
\label{subsec:lstm_ae_architecture}

The LSTM-Autoencoder (LSTM-AE) integrates the temporal modeling capacity of LSTMs with
the unsupervised representation learning of autoencoders
\cite{malhotra2015lstm, chalapathy2019deep}, yielding a model well-suited to multivariate
time-series anomaly detection. The processing pipeline is:

\begin{equation}
    X \xrightarrow{\text{Encoder LSTM}} z \xrightarrow{\text{Repeat}}
    z^{(1..T)} \xrightarrow{\text{Decoder LSTM}} H
    \xrightarrow{\text{Linear}} \hat{X}
    \label{eq:lstm_ae_pipeline}
\end{equation}

\noindent where $X \in \mathbb{R}^{T \times d}$ is the input sequence of $T$ time
steps with $d$ features, $z \in \mathbb{R}^{k}$ is the compressed latent vector, and
$\hat{X} \in \mathbb{R}^{T \times d}$ is the reconstructed sequence.

The encoder LSTM processes $X = (x_1, \ldots, x_T)$ sequentially, producing a
latent representation $z = h_T \in \mathbb{R}^{k}$ that captures the temporal
and multivariate characteristics of the window. The latent dimensionality $k$
governs the compression trade-off: too large and anomalous inputs reconstruct
faithfully; too small and normal structure is lost. The decoder repeats $z$
across the time dimension, $z^{(1..T)} = (z, \ldots, z)$, and an LSTM
reconstructs the temporal structure. A shared linear layer projects each
decoder hidden state back to the original feature space:

\begin{equation}
    \hat{x}_t = W_{\text{out}} h_t^{\text{dec}} + b_{\text{out}}
    \label{eq:linear_output}
\end{equation}

\noindent yielding $\hat{X} = (\hat{x}_1, \ldots, \hat{x}_T)$. The model is
trained end-to-end by minimizing MSE between $X$ and $\hat{X}$. The
computational footprint is governed by input dimensionality, hidden size,
number of recurrent layers, and sequence length.

\section{Bidirectional LSTM}
\label{sec:bilstm}

A Bidirectional LSTM (BiLSTM) extends the standard LSTM by processing the input
sequence in both forward ($t = 1 \rightarrow T$) and backward ($t = T \rightarrow 1$)
directions \cite{graves2005framewise, schuster1997bidirectional}, concatenating
hidden states at each step: $h_t^{\text{bi}} = [\overrightarrow{h_t} \| \overleftarrow{h_t}] \in \mathbb{R}^{2k}$.
While BiLSTMs offer richer temporal representations, they require the complete
input sequence before inference, roughly double the parameter count, and
increase memory usage. The framework uses a unidirectional LSTM-AE because edge
deployment favors lower latency and causal (past-only) inference, as discussed
in Section~\ref{sec:method_lstm_ae}.

\subsection{Model Configuration}
\label{subsec:model_config}

Table~\ref{tab:model_config} summarizes the architecture parameters used across
the experimental evaluation. The selected configuration (Config~C) balances
representational capacity with edge deployment constraints. The purpose of this
table is reproducibility: it fixes the exact model size used later in the
federated, DP, latency, and payload-size experiments. Without this table, the
reported F1 score, model footprint, and communication overhead could not be
reproduced from the narrative description alone.

\begin{table}[H]
    \centering
    \small
    \begin{tabular}{|l|c|l|}
        \hline
        \textbf{Parameter} & \textbf{Value} & \textbf{Rationale} \\
        \hline
        Input dimension ($d$) & 38 & SMD telemetry channels \\
        \hline
        Sequence length ($T$) & 50 & $\approx$8 min at 10\,s intervals \\
        \hline
        Encoder hidden ($k$) & 128 & Best F1 in sweep (Config~C) \\
        \hline
        Decoder hidden & 128 & Symmetric to encoder \\
        \hline
        LSTM layers & 2 & Depth--accuracy trade-off \\
        \hline
        Dropout & 0.2 & Regularization for edge data \\
        \hline
        Total parameters & $\approx$487\,K & Fits in $<$2\,MB serialized \\
        \hline
    \end{tabular}
    \caption{LSTM-AE architecture configuration}
    \label{tab:model_config}
\end{table}

Each row corresponds to a parameter that affects either accuracy or deployment
cost. The input dimension and sequence length define the tensor shape consumed
by the encoder. The hidden size and number of LSTM layers determine most of the
parameter count and therefore influence both reconstruction capacity and
federated payload size. Dropout controls overfitting on local edge data, while
the total parameter count confirms that the selected model remains small enough
for CPU-only edge inference.

\section{Training Procedure}
\label{sec:training_procedure}

The LSTM-AE is trained in an unsupervised fashion on sequences of normal
system behavior. The objective is to minimize the MSE loss
(Eq.~\ref{eq:mse_loss}) between input windows and their reconstructions.

\textbf{Optimization}: Adam \cite{kingma2015adam} with learning rate
$10^{-3}$, $\beta_1 = 0.9$, $\beta_2 = 0.999$. Batch size~64, trained
for 50~epochs per local round (5~epochs in FL mode).

\textbf{Threshold Calibration}: After training, reconstruction errors
are computed over the training set. The anomaly threshold is initialized
at the 95th percentile of training errors, serving as the starting point
for Q-learning adaptation (Chapter~\ref{ch:adaptive_thresholding}).

\textbf{Quantization}: Dynamic INT8 quantization is applied to LSTM and
Linear layers post-training via \texttt{torch.quantization.quantize\_dynamic},
reducing model size and inference latency without measurable accuracy
degradation.

The LSTM-AE architecture defined here provides the detection core of the
framework. Chapter~\ref{ch:framework} now embeds it within the broader
system: federated training protocol, differential privacy mechanism,
communication compression, causal root cause analysis, and adaptive
thresholding---the components that make local detection both
privacy-preserving and operationally useful.
